% Folder 79, batch 1 (archivists' pages 1-20).
% Pass: Opus 5.5 (claude-opus-5-5), 2026-10-03 — first pass, unchecked against the pages by a human.
%
% Page 1 is a cover in his hand: its words become the opening section heading
% (no \page). Page 2 is his table of contents for the whole run (author's
% pp. 1-138). Pages 3 onwards are chapter I, paginated by him from 1.
% Notation fixed for the batch: \mathfrak{S}_3 for his script S (symmetric
% group), \mathfrak{P}_k(E) for his script P (parts with k elements),
% \underline{\omega}_A for the underlined omega (set of circular orders on A),
% \widetilde{F}, \widetilde{S}_M as on the page.
\documentclass[11pt,a4paper]{article}
\input{../preamble/grothendieck}

\folder{79}
\batch{1}
\pages{1}{20}
\foldertitle{2-polyèdres réguliers triangulés, et modules libres de rangs 2 sur les anneaux $\Lambda=\mathbb{Z}/n\mathbb{Z}$ et $\Lambda=\mathbb{Z}$ : notes manuscrites (s.d.).}
\dating{[à partir de 1980]}
\watermark{Édition de démonstration}

\begin{document}

\section{2-polyèdres réguliers triangulés, et modules libres de rang 2 sur les anneaux $\Lambda=\mathbb{Z}/n\mathbb{Z}$ et $\Lambda=\mathbb{Z}$}
\note{Titre repris de la couverture (p.~1 des archivistes), de sa main ; le feuillet ne porte rien d'autre. Le titre des archivistes en dérive.}

\page{2}
\note{Table des matières de sa main ; les numéros de page à droite sont sa propre pagination de la liasse (p.~1 de l'auteur = p.~3 du fonds). Plusieurs numéros sont surchargés.}

I. Construction d'un 2-polyèdre triangulé régulier orienté… \quad p.~1

II. Relation avec l'isom. $\mathrm{Gl}(2,\mathbb{Z}) \simeq \mathfrak{G}_{\infty,3}$ \quad p.~12

III. Cas d'une structure \uncertain{½ symplectique} variable sur $M$ : recherche des structures sur $\widetilde{P}_{M} = (CS_M, F_M = \coprod_{i \in \pi} F_{M,i})$. \quad p.~18

IV. Récapitulation (1) \quad p.~30

V. Re-récapitulation, avec programme de travail \quad p.~35

VI. Calcul du \uncertain{commutant} \add{de} $\mathrm{Sl}(M)$ dans $M^*/\pm 1$, et applications \quad p.~46

Axiomatique des structures triangulées \uncertain{congruentielles} (1) \quad p.~51

VII. \uncertain{Axiomatique des} \ill{} \quad p.~65

VIII. Axiomatique : reformulation \add{\uncertain{hamiltoniens} et multigraphes} \struck{\ill{}} à \struck{\ill{}} \quad p.~81

IX. Point de vue \struck{\uncertain{graphes}} « mise en graphes » \uncertain{groupe-théorétique}… -- dualité (\ill{}) \quad p.~92

X. « Antipodismes » en général -- \uncertain{prégraphes} \uncertain{hamiltoniens} \quad \struck{\ill{}} 92

1) Antipodismes généralisés -- structures de contour \uncertain{invariante} \quad 98

2) Où les étoiles prennent le pas -- \uncertain{sur} \ill{} $\mathbb{Z}/n\mathbb{Z}$ -- \uncertain{dualité} \ill{} \quad 102

3) Présentation \uncertain{symétrique} $S, S^{\vee}$ \quad \struck{\ill{}}

4) \struck{\ill{} \uncertain{sur} condition d'\ill{}}

XI. Récapitulation des notions obtenues \quad 114

1) Graphes, prégraphes, pages, multigraphes, multipages -- \quad 120

2) Condition d'isomorphie pour un multigraphe \quad 123

3) Condition d'homogénéité. Traductions groupe-théorétiques \quad 138

4) Cas des graphes \struck{triangulants}.
\note{Les lignes VII à X sont très surchargées : les lectures y sont fragiles. La ligne non numérotée entre VI et VII porte « (1) » et la page 51.}

\subsection{I. Construction d'un 2-polyèdre triangulé régulier orienté, à un $\Lambda$-module libre de rang 2}
\note{Titre écrit en diagonale dans la marge supérieure gauche de la p.~3 ; il se termine par « p.~ex. (1) $\Lambda = \mathbb{Z}/n\mathbb{Z}$ … $\Lambda \simeq \mathbb{Z}$ », lecture incertaine.}

\page{3}
\note{p.~1 de l'auteur.}
Soit $M$ un module libre de rang 2 sur un anneau $\Lambda$. \struck{Considérons} Soit $\{x, y, z\} \in M^3$, tel que \struck{$B(M) = \{x, \ldots$}
\[
(1)\qquad x + y + z = 0
\]
alors
\[
(2)\qquad \{x,y\} \in \mathrm{Base}_\Lambda(M) \iff \{y,z\} \in \mathrm{Base}_\Lambda(M) \iff \{z,x\} \in \mathrm{Base}_\Lambda(M).
\]
On désigne par
\[
(3)\qquad B(M) \subset M^3
\]
l'ens. des triplets $(x,y,z)$ satisfaisant (1) \struck{\ill{}} et les conditions équivalentes de (2). \struck{Sous} \add{La partie} $B(M)$, \struck{\ill{}} de $M^3$ est stable par l'action de $\mathfrak{S}_3$ opérant sur $M^3$, et $\mathfrak{S}_3$ opère librement sur $B(M)$. Considérons le quotient
\[
(4)\qquad \widetilde{F}(M) \overset{\mathrm{déf}}{=} B(M)/\mathfrak{S}_3 \simeq \Bigl\{ A \subset \mathfrak{P}_3(M) \Bigm| \textstyle\sum_{x \in A} x = 0,\ \forall x \in A,\ A \setminus \{x\} \in \mathrm{Base}(M) \Bigr\}
\]
\struck{\ill{}} On peut aussi interpréter $B(M)$ comme un ens. isom. à l'ens. des \emph{repères} de $M$, i.e. des bases ordonnées $(x,y)$, par
\[
(5)\qquad B(M) \xrightarrow{\ \sim\ } \mathrm{Rep}(M) \overset{(\mathrm{déf})}{=} \mathrm{Isom}_\Lambda(\Lambda^2, M), \qquad (x,y,z) \longmapsto (x,y)
\]
l'action de $\mathfrak{S}_3$ sur $B(M)$ provient alors d'une action de $\mathfrak{S}_3$ sur $\mathrm{Rep}(M)$, déduite

\page{4}
d'une action de $\mathfrak{S}_3$ sur \add{le $\Lambda$-module} $\Lambda^2$ qu'on peut expliciter ainsi. On considère l'action de $\mathfrak{S}_3$ sur
\[
(6)\qquad \Lambda^3/\underbrace{\Lambda}_{\text{diagonale}} = V'_{\{1,2,3\}}(\Lambda) = V'_\Lambda
\]
de sorte que l'on a un isom. \struck{canonique} de $\mathfrak{S}_3$-ensembles
\[
(7)\qquad B(M) \simeq \mathrm{Isom}_\Lambda(V'_\Lambda, M)
\]
\note{Le $B$ de (7) porte au-dessus un trait surchargé (barre ou correction) ; lu ici comme le $B(M)$ de (3).}
D'ailleurs l'application (5) est composée de (7) et de la bijection
\[
\mathrm{Isom}_\Lambda(V'_\Lambda, M) \xrightarrow{\ \sim\ } \mathrm{Isom}_\Lambda(\Lambda^2, M)
\]
déduite de l'isomorphisme canonique
\[
(8)\qquad \begin{cases} \Lambda^2 \xrightarrow{\ \sim\ } \Lambda^3/\Lambda \\ (x,y) \longmapsto (x,y,0) \bmod \Lambda . \end{cases}
\]

Revenons à un $(x,y,z) \in B(M)$, notons qu'on a alors
\[
(9)\qquad x \wedge y = y \wedge z = z \wedge x \in \underbrace{\det(M)^*}_{\substack{\text{ens. des bases du } \Lambda\text{-module} \\ \text{inversible } \det M = \bigwedge^2 M}}
\]
De façon intrinsèque, si $A \in \widetilde{F}(M)$, \struck{choisissant} \add{écrivant} \struck{un ordre total} $A = \{x,y,z\}$ \add{$(x,y,z)$},

\page{5}
\note{p.~2 de l'auteur.}
l'élément $x \wedge y$ de $(\det M)^*$ ne dépend que de la permutation circulaire $\omega \in \underline{\omega}_A$ définie par l'ordre total choisi sur $A$, on trouve donc une application canonique
\[
(10)\qquad \underline{\omega}_A \xrightarrow{\ \alpha_A\ } (\det M)^*
\]
satisfaisant
\[
(11)\qquad \alpha_A(\omega^{-1}) = -\alpha_A(\omega)
\]
de sorte que $\alpha_A$ peut aussi s'interpréter comme un élément
\[
(12)\qquad \alpha_A \in \bigl(\det\nolimits_\Lambda(M) \wedge^{\mu_2} \underline{\omega}_A\bigr)^* \simeq \det\nolimits_\Lambda(M)^* \wedge^{\mu_2} \underline{\omega}_A .
\]
\struck{On} Si $M$ est « jaugé » i.e. muni d'une \emph{base} de $\det M$, \struck{soit} $e_M$, on pose
\[
(13)\qquad S\widetilde{F}(M) = \Bigl\{ A \in \widetilde{F}(M) \Bigm| \alpha_A \in \underbrace{\{\pm e_M\} \wedge^{\mu_2} \underline{\omega}_A}_{\simeq\, \underline{\omega}_A} \Bigr\}
\]
\struck{$S\widetilde{F}^+(M) = \lbrace A \in \widetilde{F}(M) \mid \alpha_A \in \ldots$}

i.e. $S\widetilde{F}(M)$ est formé des $A \in \widetilde{F}(M)$ tels que l'image de $\underline{\omega}_A$ par (10) soit $\{\pm e_M\}$. Si $2 \neq 0$ dans $\Lambda$, i.e. si $\Lambda$ n'est pas une $\mathbb{F}_2$-algèbre, alors \struck{(10) est injectif}, \struck{dès}
\[
\alpha_A : \underline{\omega}_A \xrightarrow{\ \sim\ } \{\pm e_M\}
\]
est une bijection, ce qui implique que

\page{6}
$\underline{\omega}_A$ est muni d'un élément canonique, qu'on peut noter $\omega_A$, caractérisé par
\[
(14)\qquad \begin{cases} \omega_A \in \underline{\omega}_A, \\ \alpha_A(\omega_A) = e_M \end{cases} \quad \text{ou encore} \quad \alpha_A = e_M \wedge \omega_A .
\]
Considérons aussi
\[
(15)\qquad \widetilde{S}_M = M^* \overset{\mathrm{def}}{=} \bigl\{ x \in M \bigm| x(s) \neq 0 \text{ dans } M(s) \text{ pour tout } s \in \operatorname{Spec}\Lambda \bigr\}
\]
\[
= \bigl\{ x \in M \bigm| M/\Lambda x \text{ est un } \Lambda\text{-module inversible} \bigr\}
\]
On a \struck{une application} \add{donc l'inclusion}
\[
(16)\qquad \widetilde{F}_M \subset \mathfrak{P}_3(\widetilde{S}_M = M^*)
\]
et on a en fait
\[
(17)\qquad \begin{cases} \text{si } \operatorname{Pic}(\Lambda) = 0 \quad \text{p.\,ex. si } \Lambda \text{ semi-local} \\ \text{alors } \widetilde{S}_M = M^* = \{ x \in M \mid \exists A \in \widetilde{F}_M,\ x \in A \} \\ \text{i.e. } \widetilde{S}_M = \bigcup_{A \in \widetilde{F}_M} A \end{cases}
\]
\note{« semi-local » : il écrit « ½ local ».}
Ainsi, $\widetilde{S}_M$ apparaît comme la réunion $\widetilde{F}_M$ d'un ensemble de parties à trois éléments, et on voit tout de suite que si $\{x,y\} \in \mathfrak{P}_2(\widetilde{S}_M)$, alors il existe au plus \emph{un} $A \in \widetilde{F}_M$ le contenant.

\page{7}
\note{p.~3 de l'auteur.}
Notons que $\mathrm{Gl}_\Lambda(M)$ opère sur $\widetilde{S}_M$, en stabilisant $\widetilde{F}_M$, et $\mathrm{Sl}_\Lambda(M)$ opère en stabilisant de plus le système des $\alpha_A \in \underline{\omega}_A$, $A \in \widetilde{F}_M$.

\struck{Considérons maintenant les ens.} Considérons en particulier l'action de \add{l'involution} $-\mathrm{id}_M$ sur $\widetilde{S}_M$, qui (si $2 \cdot 1_\Lambda \neq 0$) opère sans points fixes sur $\widetilde{S}_M$. On pose
\[
(18)\qquad \begin{cases} S_M = \widetilde{S}_M/\{\pm \mathrm{id}_M\} \\ F_M = \widetilde{F}_M/\{\pm \mathrm{id}_M\} . \end{cases}
\]
Notons que si $\{x,y,z\} \in \widetilde{F}_M$, alors $y \neq \pm x$, $z \neq \pm y$, $x \neq \pm z$, i.e. $\{\dot{x}, \dot{y}, \dot{z}\} \in \mathfrak{P}_3(S_M)$. On trouve ainsi une application
\[
F_M \longrightarrow \mathfrak{P}_3(S_M)
\]
je dis que celle-ci est injective, en d'autres termes, \add{si $\{x,y,z\} \in \uncertain{S_F}$ et si $\varepsilon, \varepsilon', \varepsilon'' \in \{\pm 1\}_\Lambda$} \struck{sont \ill{} des} \struck{signes} sont des signes qui ne sont pas tous égaux, alors on ne peut avoir
\[
\{\varepsilon x, \varepsilon' y, \varepsilon'' z\} \in \widetilde{F}_M .
\]
\note{Le « $S_F$ » de l'ajout interlinéaire est lu tel quel ; l'argument demande $\{x,y,z\} \in \widetilde{F}_M$.}
OPS que $\varepsilon, \varepsilon' = +1$, $\varepsilon'' = -1$, et il faut montrer

\page{8}
qu'on ne peut avoir simultanément
\[
\{x,y,z\},\ \{x,y,-z\} \in \widetilde{F}_M
\]
en effet cela implique $x + y + z = x + y - z = 0$, d'où $2z = 0$, d'où (comme $z \in M^*$) que $2 \cdot 1_\Lambda = 0$, mais alors $\varepsilon = \varepsilon' = \varepsilon''$ quand-même, contrairement à l'hypothèse.

\marginal{inutile, tout élément de $M^*$ \uncertain{fait partie} d'\ill{}}
On peut donc identifier $F_M$ à une partie de $\mathfrak{P}_3(S_M)$. Si on suppose $\operatorname{Pic}(\Lambda) = 0$, on trouve encore
\[
(19)\qquad S_M = \bigcup_{A \in \mathfrak{P}_3(S_M)} A .
\]
\note{L'indice de la réunion se lit $A \in \mathfrak{P}_3(S_M)$ ; l'argument demande $A \in F_M$, comme dans la note marginale qui suit. La note encerclée « inutile… » est reliée par un trait à « $\operatorname{Pic}(\Lambda)=0$ ».}
\marginal{si on \ill{} $\operatorname{Pic}(\Lambda) = 0$, alors (19) \ill{} $\widetilde{S}_M = S_M$ \ill{} ; il en résulte $S_M = \bigcup_{A \in F_M} A$ \ill{}}
Je désigne par $A_M$ l'ens. des parties à deux éléments \add{de $S_M$} contenues dans un \struck{$A$} $f \in F_M$ convenable
\[
(20)\qquad A_M = \{ a \in \mathfrak{P}_2(S_M) \mid \exists f \in F_M,\ a \subset f \}
\]
\[
= \{ a \in \mathfrak{P}_2(S_M) \mid \exists (x,y) \in \mathrm{Rep}(M), \text{ t.q. } \{\dot{x}, \dot{y}\} = a \}
\]
\struck{Soit} Si $a = \{\dot{x}, \dot{y}\} \in A_M$, avec $(x,y) \in \mathrm{Rep}(M)$, \add{on} cherche à déterminer l'ens. des $f \in F_M$ qui contiennent $a$. On doit donc avoir $f = \dot{A}$, où $A$ est de la forme
\note{Entre les lignes et dans la marge gauche, des ajouts très serrés, en partie illisibles : « $2 \cdot 1_\Lambda \neq 0$ », « triviales », « et $f$ unique », « On suppose \ill{} $\widetilde{F}_M \supset F_M$ », « $\widetilde{S}_M = S_M$ ». Ils ne sont pas placés.}

\page{9}
\note{p.~4 de l'auteur.}
\[
A = \{ \varepsilon x, \varepsilon' y, -(\varepsilon x + \varepsilon' y) \}, \qquad \varepsilon, \varepsilon' \in \{\pm 1_\Lambda\} .
\]
\struck{donc \ill{} est \ill{} à \ill{}} \add{et l'ens. des $f$ cherchés} s'identifie au quotient par l'involution $(\varepsilon, \varepsilon') \longmapsto (-\varepsilon, -\varepsilon')$, de l'ens. des $(\varepsilon, \varepsilon') \in \{\pm 1\}^2$. On a donc exactement deux tels $f$. En résumé :

\emph{Proposition.} Soient $\Lambda$ un anneau tel que $2 \cdot 1_\Lambda \neq 0$, et $M$ un $\Lambda$-module libre de rang 2. Soient \struck{$M^* = \{ x \in M \mid M/\Lambda x$} \struck{$\widetilde{S}_M = \{ x \in M$}
\[
(20)\qquad \begin{cases} \widetilde{S}_M = \{ x \in M \mid \exists y \in M \text{ tel que } (x,y) \in \mathrm{Rep}(M) \} \\ \text{\struck{$S_M = \widetilde{S}_M/\{\pm \mathrm{id}_M\}$}} \\ \widetilde{F}_M = \Bigl\{ \{x,y,z\} \in \mathfrak{P}_3(M) \Bigm| \begin{array}{l} x + y + z = 0 \\ (x,y) \in \mathrm{Rep}(M) \\ (\text{i.e. } (y,z) \in \mathrm{Rep}(M), \\ \text{i.e. } (z,x) \in \mathrm{Rep}(M)) \end{array} \Bigr\} \end{cases}
\]
\note{Il numérote (20) une seconde fois (cf. p.~8).}
\marginal{\ill{} définition \ill{} (15)}
\[
(21)\qquad \begin{cases} S_M = \widetilde{S}_M/\{\pm 1\} \qquad p_M : \widetilde{S}_M \to S_M \text{ appl. can.} \\ F_M = \{ f \in \mathfrak{P}_3(S_M) \mid \exists \widetilde{f} \in \widetilde{F}_M, \text{ tel que } f = p(\widetilde{f}) \} \\ A_M = \{ a \in \mathfrak{P}_2(S_M) \mid \exists f \in F_M,\ a \subset f \} \end{cases}
\]
Alors $\widetilde{F}_M$ est formé de parties à trois éléments de $\widetilde{S}_M$, de réunion $\widetilde{S}_M$, $F_M$ est formé de parties à trois éléments de $S_M$, de réunion $S_M$, et l'application \add{surj.} $\widetilde{F}_M \to F_M$, $\widetilde{f} \mapsto p(\widetilde{f})$,

\page{10}
induit une bijection canonique
\[
(22)\qquad \widetilde{F}_M/\{\pm 1\} \xrightarrow{\ \sim\ } F_M .
\]
Ceci posé, \add{pour} tout $a \in A_M$, il y a exactement deux éléments de $F_M$ qui le contiennent -- donc
\[
(23)\qquad (S_M, A_M, F_M) \overset{\mathrm{def}}{=} P_M
\]
est un 2-polyèdre combinatoire triangulé \add{$P_M$}, ayant $S_M$ comme ensemble des sommets, $A_M$ comme ens. d'arêtes, $F_M$ comme ens. des faces.
\marginal{Attention, il manque \ill{} la condition \ill{} sur l'ens. des \ill{} \ill{}}
\note{Entre « (23) » et « est un » se trouvent des ajouts interlinéaires (« \ill{} et \ill{} connexe », lecture très incertaine), reliés à la note marginale.}

\emph{Remarque.} Supposons $\Lambda$ semi-local, alors pour $x, y \in \widetilde{S}_M = M^*$, $\exists z \in \widetilde{S}_M$ tel que $(x,z) \in \mathrm{Rep}(M)$, $(y,z) \in \mathrm{Rep}(M)$, ce qui implique que pour $\dot{x}, \dot{y} \in S_M$, $\exists \dot{z} \in S_M$ tel que $\{\dot{x}, \dot{z}\} \in A_M$, $\{\dot{z}, \dot{y}\} \in A_M$. Donc la distance combinatoire de deux \struck{points} sommets de $P_M$ est au plus égale à 2, en particulier $P_M$ est un polyèdre triangulé \emph{connexe}.

\page{11}
\note{p.~5 de l'auteur.}
Soit maintenant
\[
(24)\qquad e_M \in (\det M)^*
\]
une jauge de $M$, d'où un sous-ensemble
\[
(25)\qquad S\widetilde{F}_M = \{ A \in \widetilde{F}_M \mid \alpha_A(\underline{\omega}_A) = \{\pm e_M\} \}
\]
Notons que $-\mathrm{id}_M$ applique $S\widetilde{F}_M$ dans lui-même, ce qui permet de définir
\[
(26)\qquad SF_M = S\widetilde{F}_M/\{\pm \mathrm{id}_M\} \hookrightarrow \mathfrak{P}_3(S_M)
\]
On a encore, comme on voit tout de suite
\[
(27)\qquad \widetilde{S}_M = \bigcup_{A \in S\widetilde{F}_M} A
\]
Car si $x \in \widetilde{S}_M$ \ill{} à la base $(x,y)$, alors $\exists!\, \lambda \in \Lambda^*$, \struck{\ill{} tel que} $x \wedge \lambda y = e_M$, donc $\{x, \lambda y, -(x + \lambda y)\} \in S\widetilde{F}_M$, d'où
\[
(28)\qquad S_M = \bigcup_{f \in SF_M} f .
\]
On pose
\[
(29)\qquad SA_M = \{ a \in \mathfrak{P}_2(S_M) \mid \exists f \in SF_M,\ a \subset f \}
\]
\[
= \{ a \in \mathfrak{P}_2(S_M) \mid \exists (x,y) \in \mathrm{Rep}(M),\ x \wedge y = e_M,\ a = \{\dot{x}, \dot{y}\} \}
\]
\struck{$= \lbrace a \in A_M$}

Notons que
\[
(30)\qquad SF_M \subset F_M, \quad SA_M \subset A_M
\]
sont des sous-ensembles, qui peuvent se

\page{12}
décrire ainsi. Pour $f \in F_M$, \struck{soit} on note que si $f = p(\widetilde{f})$, $\widetilde{f} \in \widetilde{F}_M$, alors \struck{$\alpha_{\widetilde{f}}(\underline{\omega}_{\widetilde{f}}) \subset (\det M)^*$ est indépendant du choix de $\widetilde{f}$, soit} l'application \add{composée}
\[
\underline{\omega}_f \xleftarrow[\text{induit par } p|\widetilde{f} \,:\, \widetilde{f} \simeq f]{} \underline{\omega}_{\widetilde{f}} \xrightarrow{\ \alpha_{\widetilde{f}}\ } (\det M)^*
\]
est indépendante de $\widetilde{f}$, soit
\[
(31)\qquad \alpha_f : \underline{\omega}_f \longrightarrow (\det M)^*
\]
et on a
\[
(32)\qquad SF_M = \{ f \in F_M \mid \alpha_f(\underline{\omega}_f) = \{\pm e_M\} \}
\]
Soit maintenant
\[
a \in \text{\struck{$S$}}A_M \qquad a = \{\dot{x}, \dot{y}\}, \text{ avec } (x,y) \in \mathrm{Rep}\, M
\]
et soient $f, f'$ les deux éléments de $F_M$ contenant $a$, donc
\[
f = \{\dot{x}, \dot{y}, (-x-y)^{\cdot}\}
\]
\[
f' = \{\dot{x}', \dot{y}' = (-y)^{\cdot}, (-x+y)^{\cdot}\}
\]
Considérons les deux applications
\note{Un premier état du diagramme, en tête, est biffé.}

(33)

\begin{tikzcd}[column sep=small, row sep=small]
 & \underline{\omega}_f \arrow[dr, "\alpha_f"] & \\
a \arrow[ur, "{\mathrm{can},\ \sim}"] \arrow[dr, "{\mathrm{can},\ \sim}"'] & & (\det M)^* \\
 & \underline{\omega}_{f'} \arrow[ur, "\alpha_{f'}"'] &
\end{tikzcd}

\page{13}
\note{p.~6 de l'auteur.}
Je dis que ce diagramme est anticommutatif, i.e. les deux applications composées $a \to (\det M)^*$ diffèrent par le signe, ou encore se déduisent l'une de l'autre par la transposition dans $a$. En effet, à l'élément $\dot{x}$, la première associe $x \wedge y$, la deuxième $x \wedge (-y)$, ce qui prouve notre assertion. Donc on a les conditions équivalentes
\[
(34)\qquad \begin{array}{l} (a \in SA_M) \iff f \in SF_M \iff f' \in SF_M \\ \qquad \iff x \wedge y \in \{\pm e_M\} \\ \qquad (\text{où } a = \{\dot{x}, \dot{y}\},\ (x,y) \in \mathrm{Rep}(M)) . \end{array}
\]
Cela montre en particulier que
\[
(35)\qquad SP_M = (S_M, SA_M, SF_M)
\]
est encore un 2-polyèdre triangulé, chaque arête étant contenue dans exactement deux faces. Il reste cependant à examiner la question de la connexité des contours combinatoires \struck{formés} associés

\page{14}
à un \struck{sommet} \add{sommet} $s \in S_M$, formé des $a \in SA_M$ et $f \in SF_M$ incidents à $s$. Choisissons un $x \in \widetilde{S}_M \subset M^*$ tel que $s = \dot{x}$, alors les $\dot{t} \in S_M$ tels que $\{s,t\} \in SA_M$ correspondent biunivoquement aux $y \in M$ tels que l'on ait
\[
(36)\qquad x \wedge y = e_M
\]
et deux tels éléments $t$ \add{$= \dot{y}$} et $u$ \add{$= \dot{z}$}, correspondant à $y, z \in M$ tels que
\[
(36')\qquad x \wedge y = x \wedge z = e_M ,
\]
sont « incidents », i.e. forment avec $s$ une face
\[
\{s,t,u\} = \{\dot{x}, \dot{y}, \dot{z}\} \in SF_M
\]
si et seulement si on peut trouver des signes $\varepsilon, \varepsilon', \varepsilon''$ tels que l'on ait
\[
\varepsilon x + \varepsilon' y + \varepsilon'' z = 0 .
\]
Quitte à multiplier par $\varepsilon$, \struck{OPS} \uncertain{on peut supposer}
\[
x + \varepsilon' y + \varepsilon'' z = 0
\]
ce qui implique, en multipliant par $x \wedge$
\[
x \wedge (\varepsilon' y + \varepsilon'' z) = 0
\]
alors qu'on a par (36')
\[
x \wedge (y - z) = 0 .
\]
Si on avait $\varepsilon' = \varepsilon''$, on aurait \struck{$x \wedge y$} aussi
\[
x \wedge (y + z) = 0
\]

\page{15}
\note{p.~7 de l'auteur.}
d'où en ajoutant
\[
2\, x \wedge y = 0 \quad \text{i.e.} \quad 2 e_M = 0
\]
ce qui impliquerait $2 \cdot 1_\Lambda = 0$, contrairement à l'hypothèse. On doit donc avoir
\[
(37)\qquad \begin{cases} x + y - z = 0 \quad \text{ou} \quad x - y + z = 0 \\ \text{i.e.} \quad z = x + y \quad \text{ou} \quad y = x + z \end{cases}
\]
Donc, l'ens. des sommets de $SP_M$ incidents à $s \in S_M$, $s = \dot{x}$, s'identifie à la droite affine $L_x$ sur $\Lambda$ formée des $y \in M$ satisfaisant (36), avec la structure de contour où les « arêtes »
\[
\{y,z\} \in \mathfrak{P}_2(L_x) \qquad L_x = \{ y \in M \mid x \wedge y = e_M \}
\]
\struck{sont \uncertain{les}} du contour sont formées des paires $\{y,z\}$ telles que l'on ait
\[
y - z \in \{\pm x\} .
\]
Les composantes connexes des contours correspondent \add{donc} aux \struck{\ill{}} \add{les} parties de $L_x$ de la forme
\[
(38)\qquad \{ y + n x \mid n \in \mathbb{Z} \} = y + \Lambda_0 x
\]

\page{16}
où
\[
(39)\qquad \Lambda_0 = \operatorname{Im}(\mathbb{Z} \to \Lambda), \quad n \mapsto n \cdot 1_\Lambda
\]
donc l'ens. des composantes connexes est l'ens. \struck{des} quotient
\[
(40)\qquad L_x/\Lambda_0 = L_x \wedge_\Lambda \Lambda/\Lambda_0 .
\]
Ainsi, la structure $SP_M$ est une vraie structure \add{de 2-polyèdre} triangulé, \struck{\ill{}} i.e. satisfait à l'axiome de connexité pour le contour associé à un sommet, ssi on a $\Lambda_0 = \Lambda$, i.e. ssi on a
\[
(41)\qquad \Lambda \simeq \mathbb{Z} \quad \text{ou} \quad \Lambda \simeq \mathbb{Z}/n\mathbb{Z} \quad \text{pour } n \in \mathbb{N}^*,\ n \geq 3
\]
(NB on a exclu le cas $n = 1$ ou $2$, avec la condition $2 \cdot 1_\Lambda \neq 0$.)
\note{« $L_x \wedge_\Lambda \Lambda/\Lambda_0$ » est lu tel quel ; on attend $L_x \otimes_\Lambda \Lambda/\Lambda_0$.}

\emph{Remarque.} Si on avait travaillé dans $P_M$ au lieu de $SP_M$, on trouverait \struck{pour l'ens. des sommets liés à} \struck{$s = \dot{x}$} \struck{l'ens. des sommets de $P_M$ est en corresp. \ill{} avec l'ens. quotient} \add{l'ens. des sommets liés à $s = \dot{x}$ … } … et la connexité ssi on … \ill{}

\note{La Remarque de la p.~16 est très raturée ; les ajouts interlinéaires ne se raccordent pas avec certitude. La fin du feuillet se lit « … connexité ssi on … », la suite à la page suivante.}

\page{17}
\note{p.~8 de l'auteur.}
a $SP_M = P_M$ i.e. $SF_M = F_M$ i.e. $\Lambda^* \simeq \{\pm 1\}$, \emph{et} $\Lambda = \Lambda_0$ i.e. $\Lambda \simeq \mathbb{Z}$ ou $\Lambda \simeq \mathbb{Z}/n\mathbb{Z}$, $n \geq 3$. Notons que la structure connue de $(\mathbb{Z}/n\mathbb{Z})^*$ montre alors que les seuls cas sont ceux \struck{\ill{}}
\[
(42)\qquad \Lambda = \mathbb{Z}, \quad \Lambda = \mathbb{Z}/n\mathbb{Z} \text{ avec } n \in \{3,4,6\} .
\]

Revenons au cas général, \struck{et regardant} et notons que $S\widetilde{F}_M$, donc aussi $SP_M$, ne \struck{dépend} change pas si on remplace la jauge $e_M$ par $-e_M$. Ainsi le groupe
\[
(43)\qquad \mathrm{Gl}_\Lambda(M)^{\pm} \overset{\mathrm{def}}{=} \{ g \in \mathrm{Gl}_\Lambda(M) \mid \det g \in \{\pm 1\} \}
\]
opère sur le 2-polyèdre (généralisé) $SF_M$. Nous allons introduire sur celui-ci une \emph{orientation}, qui elle dépend du choix de $e_M$, et est remplacée par l'orientation opposée, quand on remplace $e_M$ par $-e_M$. Pour ceci, revenons à l'application (31)
\[
\alpha_f : \underline{\omega}_f \longrightarrow (\det M)^*
\]

\page{18}
associée à un $f \in F_M$, qui pour $f \in SF_M$ devient une bijection
\[
(43)\qquad \alpha_f : \underline{\omega}_f \xrightarrow{\ \sim\ } \{\pm e_M\}
\]
\note{Il numérote (43) une seconde fois (cf. p.~17).}
d'où, grâce au choix de $e_M$, un élément
\[
(44)\qquad \omega_f \in \underline{\omega}_f \qquad \alpha_f(\omega_f) = e_M ,
\]
\struck{\ill{}} i.e. une \emph{orientation} \add{\emph{canonique}} de la face $f$. La propriété de compatibilité (33) montre que les orientations des faces \add{adjacentes} sont compatibles deux à deux, \add{donc} définissent une orientation \add{$\mu$} de $SP_M$, appelée \emph{l'orientation canonique}. On voit donc que le groupe $\mathrm{Gl}_\Lambda(M)^{\pm}$ opère sur $SP_M$, et opère sur l'ens. des orientations \struck{\ill{} en échangeant les \ill{}} \add{l'orientation canonique} suivant $\det(g)$ :
\[
(45)\qquad g \cdot \mu = \text{\struck{$\ldots$}}\ \det(g) \cdot \mu \qquad g \in \mathrm{Gl}_\Lambda(M)^{\pm 1}
\]
Revenons \add{d'abord} à la question de la connexité

\page{19}
\note{p.~9 de l'auteur.}
de $SF_M$, donc à la question, pour deux éléments $x, y \in \widetilde{S}_M$, de trouver une suite d'éléments
\[
(46)\qquad \begin{cases} z_0, z_1, \ldots, z_N \in M \\ z_0 = x,\ z_N = y, \qquad z_i \wedge z_{i+1} \in \{\pm e_M\} \text{ pour } 0 \leq i \leq N-1 \end{cases}
\]
Supposons d'abord que $(x,y) \in \mathrm{Rep}(M)$, alors on aura
\[
e_M = \lambda\, x \wedge y \qquad \lambda \in \Lambda^*
\]
soit donc \struck{$z = \ldots$}
\[
z = \lambda x + \lambda y
\]
alors
\[
x \wedge z = z \wedge y = e_M
\]
et on gagne, en prenant la chaîne $z_0 = x$, $z_1 = z$, $z_2 = y$. Cela nous ramène donc à relier $x$ à $y$ par une chaîne d'éléments $z_i$, tels que deux consécutifs forment une base, ce qui était toujours possible si $\Lambda$ était semi-local -- c'est donc OK pour le cas des anneaux $\mathbb{Z}/n\mathbb{Z}$, \struck{\ill{}} avec $n \in \mathbb{N}^*$. \struck{Mais on aimerait aussi -- } \marginal{On trouve}
\note{Le « $z = \lambda x + \lambda y$ » : le signe entre les deux termes est surchargé.}

\page{20}
s'assurer de la connexité dans le cas $\Lambda = \mathbb{Z}$.

Un complément dans le cas semi-local : \struck{On} sait que $\forall x, y \in M^*$, $\exists z \in M^*$, tels que
\[
(x,z) \in \mathrm{Rep}(M), \quad (z,y) \in \mathrm{Rep}(M),
\]
\struck{On} aura donc
\[
x \wedge z = \lambda e_M, \quad z \wedge y = \mu e_M \qquad \lambda, \mu \in \Lambda^*
\]
et quitte à remplacer $z$ par $\frac{1}{\lambda} z$, \struck{OPS}
\[
x \wedge z = e_M, \quad z \wedge y = \mu e_M
\]
et par l'argument précédent on trouve $z' \in M$ tel qu'on ait aussi $z \wedge z' = z' \wedge y = e_M$, donc
\[
(47)\qquad x \wedge z = z \wedge z' = z' \wedge y = e_M
\]
ce qui montre que deux sommets de $SF_M$ (ou de $F_M$) sont à une distance combinatoire $\leq 3$. (NB. La borne est atteinte pour l'icosaèdre, correspondant au cas de $\Lambda = \mathbb{Z}/5\mathbb{Z}$.)

Pour le cas de $\Lambda = \mathbb{Z}$, particulièrement intéressant du point de vue « univers », il est
\note{La phrase se poursuit au-delà du lot (p.~21 des archivistes, p.~10 de l'auteur).}

\end{document}
